\documentclass[11pt]{article}
\usepackage[right=1in,left=1in,top=0.85in,bottom=0.85in]{geometry}
\usepackage{amsmath,amssymb,amsfonts}
\usepackage{booktabs}
\usepackage{algorithm}
\usepackage{algorithmic}
\usepackage{hyperref}
\usepackage{placeins}
\usepackage{natbib}
\usepackage{titling}
\setlength{\droptitle}{-2.5em}
\posttitle{\par\end{center}\vspace{-1.2em}}
\setlength{\parskip}{0.3em}
\bibliographystyle{abbrvnat}

\newcommand{\E}{\mathbb{E}}
\newcommand{\R}{\mathbb{R}}
\newcommand{\calA}{\mathcal{A}}
\newcommand{\calS}{\mathcal{S}}
\newcommand{\calH}{\mathcal{H}}
\newcommand{\calD}{\mathcal{D}}

\title{Primer: Adversarial Inverse Reinforcement Learning (AIRL)\\[3pt]
\large Fu, Luo, and Levine (2018)}
\author{econirl package documentation}
\date{}

\begin{document}
\maketitle
\vspace{-0.5em}

%% ============================================================
\section{Context and Problem}
%% ============================================================

Maximum entropy IRL \citep{ziebart2010} recovers a reward function that rationalizes observed behavior, but the recovered reward entangles the true reward with the value function.
For any bounded $h : \calS \to \R$, the shaped reward
\begin{equation}
\label{eq:shaping}
r'(s,a) \;=\; r(s,a) \;+\; \gamma\, \E_{s' \sim P(\cdot|s,a)}[h(s')] \;-\; h(s)
\end{equation}
produces the same optimal policy under the original transitions $P$ \citep{ng1999}.
If $P$ changes (new environment, counterfactual), the shaping term evaluates differently and the shaped reward no longer rationalizes the same behavior.
\citet{fu2018} propose AIRL to disentangle the transferable reward from the non-transferable value function by structuring the discriminator so that one component converges to $r^*(s)$ and the other absorbs $V^*(s)$.

%% ============================================================
\section{Model and Notation}
%% ============================================================

The agent maximizes entropy-regularized cumulative reward \citep{ziebart2010}:
\begin{equation}
\label{eq:objective}
\max_\pi \; \E_\pi\!\left[\sum_{t=0}^{\infty} \gamma^t \big(r(s_t, a_t) + \calH(\pi(\cdot \mid s_t))\big)\right].
\end{equation}
The optimal policy satisfies $\pi^*(a \mid s) = \exp(A^*) / \sum_{a'} \exp(A^*_{a'})$ where $A^*(s,a) = Q^*(s,a) - V^*(s)$, giving $\log \pi^*(a \mid s) = A^*(s,a)$.

\begin{table}[h!]
\centering\small
\begin{tabular}{ll}
\toprule
Symbol & Meaning \\
\midrule
$s \in \calS$, $a \in \calA$, $s'$ & State, action, next state \\
$\gamma$, $P(s'|s,a)$ & Discount factor, transition kernel \\
$r^*(s)$ & True reward (state-only; required for Theorem~1) \\
$\pi^*(a|s)$, $V^*$, $Q^*$, $A^*$ & Expert policy, value, action-value, advantage \\
\midrule
$g_\theta(s)$ & Reward approximator (state-only, param.\ $\theta$) \\
$h_\phi(s)$ & Shaping potential $\approx V(s)$ (param.\ $\phi$) \\
$\pi_\omega(a|s)$ & Learned policy (param.\ $\omega$) \\
$f_{\theta,\phi}(s,a,s')$ & Discriminator logit \eqref{eq:f}; \quad $D_{\theta,\phi} \in [0,1]$: discriminator output \eqref{eq:D} \\
\bottomrule
\end{tabular}
\end{table}
\FloatBarrier

AIRL decomposes the discriminator into reward and shaping:
\begin{align}
f_{\theta,\phi}(s,a,s') &= g_\theta(s) + \gamma\, h_\phi(s') - h_\phi(s), \label{eq:f}\\[2pt]
D_{\theta,\phi}(s,a,s') &= \frac{\exp(f_{\theta,\phi})}{\exp(f_{\theta,\phi}) + \pi_\omega(a \mid s)}. \label{eq:D}
\end{align}
The logit $f$ mirrors the Bellman relationship $Q = r + \gamma \E[V(s')]$, decomposed into $g_\theta \approx r$ and $h_\phi \approx V$.
Note $g_\theta$ depends on $s$ only, not on $a$.

%% ============================================================
\section{Identification}
%% ============================================================

\paragraph{Reward ambiguity.}
IRL rewards are identified only up to potential-based shaping~\eqref{eq:shaping} \citep{ng1999}.
The observed policy constrains $|\calS|(|\calA|-1)$ values but the reward has $|\calS| \cdot |\calA|$ unknowns, leaving $|\calS|$ free variables corresponding to the value function \citep{cao2021}.

\paragraph{Decomposability (Definition B.1).}
An MDP is \emph{decomposable} if all states are linked: for every pair $(s, s')$ there exists a policy and time $t$ such that $s'$ is reachable from $s$ with positive probability.
This rules out absorbing states.

\paragraph{Disentanglement (Theorem 1).}
If (a)~$g_\theta$ is state-only, (b)~the MDP is decomposable, and (c)~the adversarial training converges to the Nash equilibrium, then $g_\theta(s) = r^*(s) + c$ for a constant $c$.

\emph{Proof sketch.}
At Nash equilibrium $D = 1/2$, so $f = \log \pi^* = A^*$.
Substituting~\eqref{eq:f}: $g_\theta(s) + \gamma h_\phi(s') - h_\phi(s) = r^*(s) + \gamma V^*(s') - V^*(s)$.
Since $g_\theta$ is state-only, taking the expectation over $s'$ yields $g_\theta(s) = r^*(s) + h_\phi(s) - V^*(s) + c_0$.
Decomposability implies $h_\phi(s) - V^*(s)$ is constant, so $g_\theta(s) = r^*(s) + c$. \hfill $\square$

\paragraph{Limitations.}
Theorem~1 requires state-only rewards and excludes absorbing states.
DDC models have action-dependent payoffs and exit options.
\citet{lee2026} resolve both with anchor constraints (see AAIRL primer).

%% ============================================================
\section{Estimation Algorithm}
%% ============================================================

\begin{algorithm}[h!]
\caption{AIRL (Fu et al.\ 2018, Algorithm 1)}
\label{alg:airl}
\begin{algorithmic}[1]
\REQUIRE Expert demonstrations $\calD_E = \{\tau_1, \ldots, \tau_N\}$, discount $\gamma$
\ENSURE Disentangled reward $g_\theta(s)$

\STATE Initialize $\theta_0, \phi_0, \omega_0$

\FOR{iteration $= 1, 2, \ldots$}

\STATE $\calD_\pi \leftarrow$ roll out $\pi_\omega$ to collect $M$ transitions \hfill \textit{// sample from policy}

\FOR{$j = 1, \ldots, J$} \hfill \textit{// train discriminator}
  \STATE $f_n \leftarrow g_\theta(s_n) + \gamma\, h_\phi(s'_n) - h_\phi(s_n)$ \hfill \textit{// Eq.~\eqref{eq:f}}
  \STATE $D_n \leftarrow \sigma(f_n - \log \pi_\omega(a_n \mid s_n))$ \hfill \textit{// Eq.~\eqref{eq:D}}
  \STATE $\calL \leftarrow -\E_{\calD_E}[\log D] - \E_{\calD_\pi}[\log(1-D)]$
  \STATE $\theta, \phi \leftarrow \theta, \phi - \alpha_D \nabla_{\theta,\phi} \calL$
\ENDFOR

\STATE $\tilde{r} \leftarrow \log D - \log(1-D) = f - \log \pi_\omega$ \hfill \textit{// reward signal}
\STATE Update $\omega$ via PPO on $\max_{\pi_\omega} \E[\sum_t \gamma^t \tilde{r}]$ \hfill \textit{// policy update}

\IF{$D \approx 1/2$ on held-out data}
  \STATE \textbf{break}
\ENDIF

\ENDFOR
\RETURN $g_\theta$ \quad {\small (disentangled, transferable reward)}
\end{algorithmic}
\end{algorithm}
\FloatBarrier

\citet{fu2018} use PPO \citep{schulman2017} with $J = 1$ discriminator step per policy step.
The reward signal $\tilde{r} = f - \log \pi_\omega$ is the log-odds of the discriminator.
At convergence $D = 1/2$ and $g_\theta = r^* + c$ by Theorem~1.

%% ============================================================
\section{Results from the Paper}
%% ============================================================

\paragraph{Transfer experiment (Table 1 in paper).}
AIRL is trained on a PointMaze where the agent navigates to a goal in a room with one wall layout.
The recovered reward $g_\theta$ is then used to train a new policy in a maze with different walls.
AIRL transfers successfully: the agent reaches the goal in the new maze because $g_\theta$ captures the goal location, not the path around specific walls.
GAIL and MaxEnt IRL fail transfer because their recovered rewards encode the original wall layout through the entangled shaping potential.
This is the paper's central empirical result.

\paragraph{Continuous control (Table 2).}
On Ant, HalfCheetah, and Hopper from OpenAI Gym, AIRL achieves cumulative reward within 5\% of the expert on all three tasks.
GAIL matches this performance but does not recover a reward function.
MaxEntIRL underperforms on the higher-dimensional tasks (Ant) because its linear reward parameterization is too restrictive for continuous control.

\paragraph{State-only vs.\ state-action (Section 5.2).}
When the true reward is state-only, $g_\theta(s)$ recovers it up to a constant as Theorem~1 predicts.
When the true reward depends on actions, the state-only $g_\theta(s)$ absorbs the action-dependent component into the shaping potential $h_\phi$.
The reward is still useful for the original environment but does not transfer cleanly because the action-dependent piece is entangled with the dynamics.
This motivates the AAIRL extension by \citet{lee2026}, which handles action-dependent rewards via anchor constraints.

%% ============================================================
\section{Simulation}
%% ============================================================

We test on a $5 \times 5$ gridworld: 25 states, 4 actions (N/S/E/W), stochastic transitions (80\% intended, 20\% random), $\gamma = 0.95$.
True reward: $r^*(s) = 1$ at goal (bottom-right), 0 elsewhere.
Expert data: 500 trajectories of 50 steps from the optimal soft Bellman policy.
Run \texttt{python airl\_run.py} to regenerate Table~\ref{tab:results}.

\begin{table}[h!]
\centering\small
\begin{tabular}{lrr}
\toprule
Metric & True & AIRL \\
\midrule
% Auto-generated by airl_run.py -- do not edit by hand
Cosine similarity & $---$ & $0.000$ \\
Correlation & $---$ & $0.000$ \\
Policy match rate & $---$ & $0.000$ \\
Goal reward & $1.00$ & $0.000$ \\
Mean non-goal reward & $0.00$ & $0.000$ \\

\bottomrule
\end{tabular}
\caption{Reward recovery on 25-state gridworld (state-only reward, decomposable MDP). Auto-generated by \texttt{airl\_run.py}.}
\label{tab:results}
\end{table}
\FloatBarrier

The decomposability condition holds because the 20\% slip probability makes all states reachable from all others.
AIRL with conservative policy iteration ($\alpha_p = 0.1$) and weight decay should recover the goal reward up to a constant.
Cosine similarity measures how well the recovered reward vector aligns with the true reward vector in direction (ignoring scale).
Policy match rate measures the fraction of states where the greedy action under the recovered reward agrees with the expert's greedy action.
The econirl implementation uses \texttt{AIRLEstimator} with \texttt{AIRLConfig(reward\_type="tabular", policy\_step\_size=0.1, use\_shaping=True)}.

%% ============================================================
\section{Comparison}
%% ============================================================

\begin{table}[h!]
\centering\small
\begin{tabular}{lccccc}
\toprule
 & MaxEnt & MCE-IRL & GAIL & AIRL & AAIRL \\
\midrule
Recovers reward & Yes & Yes & No & Yes & Yes \\
Reward transfers & No & No & N/A & \textbf{Yes} & Yes \\
Action-dep.\ $r$ & Yes & Yes & N/A & No$^*$ & Yes \\
Continuous $\calS$ & No & No & Yes & Yes & Yes \\
Requires $P$ & Yes & Yes & No & No & No$^\dagger$ \\
\bottomrule
\end{tabular}
\caption{$^*$Theorem~1 requires state-only $g_\theta(s)$. $^\dagger$With PPO generator; tabular variant uses known $P$.}
\label{tab:comparison}
\end{table}
\FloatBarrier

AIRL's contribution over MaxEnt/MCE-IRL is reward transferability via disentanglement.
The shaped reward recovered by MaxEnt IRL performs well in the training environment but fails when dynamics change, because the shaping potential $h$ is calibrated to the original transitions.
AIRL avoids this by separating $g_\theta$ (transferable) from $h_\phi$ (non-transferable) within the discriminator architecture.

Over GAIL, AIRL recovers a reward function, not just a policy.
GAIL trains a discriminator that classifies expert vs.\ generated transitions, but the discriminator's internal representation is unstructured.
AIRL imposes the Bellman-consistent decomposition~\eqref{eq:f}, which gives the discriminator components interpretable meaning.

The main limitation is the state-only restriction.
Theorem~1 requires $g_\theta(s)$ to depend on $s$ only, not on $a$.
In DDC models, different actions carry different payoffs (buying costs money, waiting costs time, exiting is free), so the reward is inherently action-dependent.
A state-only $g_\theta$ cannot capture these differences and absorbs them into the shaping potential, breaking the disentanglement guarantee.
\citet{lee2026} resolve this with anchor constraints in AAIRL: setting $r(s, a_{\text{exit}}) = 0$ and $V(s_{\text{abs}}) = 0$ pins down the action-dependent reward uniquely, extending AIRL to the full DDC setting.

\bibliography{references}

\end{document}
